% AP.D — Golden Mirror (Trinity ↔ Flos Aureus symmetry)
% φ² + φ⁻² = 3 · Trinity S³AI
% Author: Dmitrii Vasilev (ORCID 0009-0008-4294-6159)
% Issue: #380 task 2.4

\chapter*{Appendix D: Golden Mirror — Trinity \texorpdfstring{$\leftrightarrow$}{<->} Flos Aureus Symmetry}
\addcontentsline{toc}{chapter}{Appendix D: Golden Mirror}
\label{app:golden-mirror}

\textbf{Anchor invariant:} \(\varphi^2 + \varphi^{-2} = 3\).

\section*{D.1 Two-Pillar Structure}

The dissertation is organised around two complementary pillars, united by
the Trinity anchor \(\varphi^2 + \varphi^{-2} = 3\):

\begin{itemize}
  \item \textbf{Trinity S\textsuperscript{3}AI} (35 chapters, Ch.0..34):
    hardware-numerics, empirical AI benchmarks, FPGA coprocessor
  \item \textbf{Flos Aureus} (34 chapters, FA.00..33):
    sacred geometry, algebraic foundations, mathematical beauty
\end{itemize}

This appendix documents the structural symmetry between the two pillars
and proves that both derive from the same algebraic root.

\section*{D.2 Symmetry Table}

\begin{longtable}{lll}
\toprule
\textbf{Trinity S\textsuperscript{3}AI concept} &
\textbf{Mirror (Flos Aureus)} &
\textbf{Shared root} \\
\midrule
GoldenFloat GF16 & FA.06 Golden Mantissa & \(\varphi^2 = \varphi + 1\) \\
Period-Locked Monitor & FA.17 Golden Spiral & \(\varphi\)-continued fraction \\
VSA XOR bind & FA.12 Flower of Life & Binary $\oplus$ = tiling \\
Fibonacci dimension seeds & FA.07 Golden Sprout & \(F_n\) recursion \\
TF3 ternary format & FA.01 Golden Egg & \{$-1,0,+1$\} = 3 states \\
KOSCHEI ISA & FA.14 Platonic Solids & Icosahedral symmetry \\
TRI27 DSL & FA.13 Metatron's Cube & 13-node \(\varphi\)-graph \\
GF4..GF64 family & FA.23 GF(16) Algebra & Galois field $\subset \varphi$ \\
BPB benchmark & FA.25 Benchmarks & Information entropy \\
Energy 3000$\times$ & FA.09 Golden Seal & Minimal surface area \\
R-rules (constitutional) & FA.32 Conclusion & Invariant preservation \\
\(\varphi^2 + \varphi^{-2} = 3\) & \(\varphi^2 + \varphi^{-2} = 3\) & \textbf{THE ANCHOR} \\
\bottomrule
\caption{Golden Mirror correspondence table. Every Trinity S\textsuperscript{3}AI
concept maps to a Flos Aureus chapter via the shared \(\varphi\)-root.}
\end{longtable}

\section*{D.3 Proof of Shared Root}

\begin{theorem}[Golden Mirror Duality — THM-D.1]
\label{thm:D:1}
For every algebraic identity \(I(\varphi)\) used in Trinity S\textsuperscript{3}AI
(hardware layer), there exists a dual geometric identity \(I'(\varphi)\)
in Flos Aureus (geometry layer) such that \(I(\varphi) = I'(\varphi)\) as
polynomials over \(\mathbb{Q}(\varphi)\).
\end{theorem}

\begin{proof}
Both pillars ground identities in \(\varphi^2 = \varphi + 1\), the minimal
polynomial of the golden ratio over \(\mathbb{Q}\). Every identity
\(I(\varphi)\) in either pillar reduces to a polynomial in \(\{1, \varphi\}\)
(since \(\mathbb{Q}(\varphi)\) is 2-dimensional over \(\mathbb{Q}\)).
The reduction algorithm is the Euclidean division by \(X^2 - X - 1\).
Therefore \(I\) and \(I'\) share coefficients in \(\mathbb{Q}\), proving
equality as polynomials.
\end{proof}

\section*{D.4 Chapter-Count Symmetry}

\begin{table}[h]
\centering
\begin{tabular}{lrr}
\toprule
\textbf{Pillar} & \textbf{Chapters} & \textbf{Body chars} \\
\midrule
Trinity S\textsuperscript{3}AI (Ch.0..34) & 35 & $\approx$590,000 \\
Flos Aureus (FA.00..33) & 34 & $\approx$1,940,000 \\
\midrule
Total & 69 + 29 front+appendix & 2,530,000 \\
\bottomrule
\end{tabular}
\caption{The ratio $1{,}940 / 590 \approx 3.29 \approx 2\varphi$ (within 1.9\%).}
\end{table}

The character ratio of the two pillars approximates \(2\varphi \approx
3.236\). This is not a design constraint but an emergent property of
writing Flos Aureus chapters to full geometric depth while Trinity chapters
are bounded by experimental data.

\(\varphi^2 + \varphi^{-2} = 3\)
